\documentclass[a4paper,12pt]{article}
% Use native Unicode fonts with XeLaTeX/LuaLaTeX; retain T5 for pdfLaTeX.
\usepackage{iftex}
\ifPDFTeX
  \usepackage[utf8]{vntex}
\else
  \usepackage{fontspec}
  \defaultfontfeatures{Ligatures=TeX}
  \IfFontExistsTF{Times New Roman}
    {\setmainfont{Times New Roman}}
    {\setmainfont{TeX Gyre Termes}}
  \renewcommand{\contentsname}{Mục lục}
  \renewcommand{\refname}{Tài liệu tham khảo}
  \renewcommand{\figurename}{Hình}
  \renewcommand{\tablename}{Bảng}
\fi
\usepackage[margin=2.3cm,headheight=36pt]{geometry}
\usepackage{graphicx,longtable,booktabs,array,caption,float,fancyhdr,lastpage,xcolor,amsmath}
\usepackage[unicode,colorlinks=true,linkcolor=black,urlcolor=blue]{hyperref}
\usepackage{xurl}
\usepackage{enumitem}
\setlist{nosep,leftmargin=*}
\setlength{\parindent}{0pt}
\setlength{\parskip}{6pt}
\setlength{\emergencystretch}{3em}
\renewcommand{\arraystretch}{1.25}
\captionsetup{font=small,labelfont=bf}
\newcolumntype{L}[1]{>{\raggedright\arraybackslash}p{#1}}
\pagestyle{fancy}
\fancyhf{}
\fancyhead[L]{\small\textbf{TRƯỜNG ĐẠI HỌC SÀI GÒN}\\Khoa Công nghệ Thông tin}
\fancyhead[R]{\small BrewLite -- Review 1}
\fancyfoot[L]{\footnotesize Công nghệ phần mềm -- Học kỳ I, 2026--2027}
\fancyfoot[R]{\footnotesize Trang \thepage/\pageref{LastPage}}
% Data the team still has to supply.
\newcommand{\todo}[1]{\textcolor{red!65!black}{[Cần bổ sung: #1]}}
% Screenshot slot: shows figures/<file> when present, otherwise a framed placeholder.
\newcommand{\khunganh}[2]{%
  \IfFileExists{figures/#2}%
    {\includegraphics[width=#1]{figures/#2}}%
    {\fbox{\parbox[c][3cm][c]{\dimexpr#1-2\fboxsep-2\fboxrule\relax}{\centering\footnotesize Chưa có ảnh\\\texttt{figures/#2}}}}}
% Postman case table: code, what is sent, expected result.
\newenvironment{bangca}{%
  \begin{longtable}{L{0.11\linewidth}L{0.47\linewidth}L{0.32\linewidth}}
  \toprule Mã ca & Kiểm tra & Kết quả mong đợi\\\midrule\endfirsthead
  \toprule Mã ca & Kiểm tra & Kết quả mong đợi\\\midrule\endhead
  \bottomrule\endfoot}
  {\end{longtable}}
\begin{document}
\begin{titlepage}
\centering
{\large TRƯỜNG ĐẠI HỌC SÀI GÒN\\KHOA CÔNG NGHỆ THÔNG TIN\par}
\vspace{1cm}
\includegraphics[width=3cm]{logoITSGU.png}
\vspace{1cm}
\hrule\vspace{0.5cm}
{\Large\bfseries BÁO CÁO BÀI TẬP LỚN\\CÔNG NGHỆ PHẦN MỀM\par}
\vspace{0.6cm}
{\Huge\bfseries BrewLite\par}
\vspace{0.4cm}
{\large Ứng dụng đặt cà phê không dùng tiền mặt\\Triển khai 10 task theo Agile Scrum\par}
\vspace{0.5cm}\hrule
\vspace{0.8cm}
{\large BẢN DEMO -- REVIEW LẦN 1\par}
\vspace{0.5cm}
\begin{tabular}{lll}
\toprule Thành viên & MSSV & Phân công\\\midrule
Bảo Tín & \todo{MSSV} & Task 1, 6\\
Phú & \todo{MSSV} & Task 2, 9\\
Bảo Thiện & \todo{MSSV} & Task 3, 5\\
Lợi & \todo{MSSV} & Task 4, 7\\
Thuận & \todo{MSSV} & Task 8, 10\\\bottomrule
\end{tabular}
\vfill
GVHD: Từ Lãng Phiêu\\
Học kỳ I -- Năm học 2026--2027\\
TP. Hồ Chí Minh, tháng 10 năm 2026
\end{titlepage}
\tableofcontents\newpage

\section{Tóm tắt và phạm vi}
BrewLite cho khách chọn đồ uống, chọn size và topping, quản lý giỏ hàng, đăng nhập, tạo đơn, thanh toán bằng Ví hoặc Thẻ qua cổng giả lập và xem lịch sử đơn. Frontend viết bằng Next.js, backend bằng NestJS, dữ liệu lưu trong PostgreSQL qua Prisma. Docker Compose dựng ba dịch vụ: cơ sở dữ liệu ở cổng 5432, backend ở cổng 3000, frontend ở cổng 3001.

Đề chia sản phẩm thành 10 task, mỗi task 1 điểm. Báo cáo đối chiếu từng task với tiêu chí ở mục 8 và mục 9 của đề. Số liệu kiểm thử lấy từ lần chạy ngày 21/10/2026 trên nhánh \texttt{develop-agile}.

\section{Yêu cầu và kiến trúc}
\subsection{Bên liên quan}
\begingroup\small
\begin{longtable}{L{0.22\linewidth}L{0.48\linewidth}L{0.22\linewidth}}
\toprule Bên liên quan & Việc làm được trên BrewLite & Kênh\\\midrule\endfirsthead
\toprule Bên liên quan & Việc làm được trên BrewLite & Kênh\\\midrule\endhead
\bottomrule\endfoot
Khách hàng & Xem menu, chọn size và topping, đặt đơn, thanh toán, hủy đơn của mình, xem lịch sử và điểm thưởng & Web\\
Nhân viên quầy (Barista) & Xem hàng đợi đơn đã thanh toán, chuyển đơn sang các bước pha chế, hủy đơn đã thanh toán khi hết nguyên liệu & REST API, kiểm bằng Postman\\
Quản lý (Admin) & Mọi việc của Barista, xem kho kể cả món hết hàng, chỉnh tồn kho & REST API, kiểm bằng Postman\\
Cổng thanh toán giả lập & Chấp nhận hoặc từ chối giao dịch Ví/Thẻ & Backend gọi trực tiếp\\
\end{longtable}
\endgroup
Web không có màn hình cho Barista và Admin. Mục 3.2 của đề xếp chức năng quầy ngoài MVP, nên nhóm chỉ làm phần API phục vụ máy trạng thái đơn hàng của Task 10.

\subsection{Yêu cầu chức năng và phi chức năng}
\begingroup\small
\begin{longtable}{L{0.08\linewidth}L{0.62\linewidth}L{0.22\linewidth}}
\toprule STT & Yêu cầu chức năng (mục 4.1 của đề) & Task\\\midrule\endfirsthead
\toprule STT & Yêu cầu chức năng (mục 4.1 của đề) & Task\\\midrule\endhead
\bottomrule\endfoot
1 & Hiển thị menu đồ uống kèm hình, tên, giá & Task 2, 3\\
2 & Xem chi tiết sản phẩm, chọn size và topping & Task 4\\
3 & Thêm, sửa, xóa sản phẩm trong giỏ; tự tính tổng tiền & Task 5\\
4 & Đăng ký, đăng nhập bằng JWT & Task 7\\
5 & Tạo đơn và thanh toán không tiền mặt & Task 6, 8\\
6 & Hiển thị xác nhận đơn và lịch sử đơn & Task 9\\
\end{longtable}

\begin{longtable}{L{0.30\linewidth}L{0.62\linewidth}}
\toprule Yêu cầu phi chức năng (mục 4.2 của đề) & Minh chứng\\\midrule\endfirsthead
\toprule Yêu cầu phi chức năng (mục 4.2 của đề) & Minh chứng\\\midrule\endhead
\bottomrule\endfoot
API phản hồi dưới 500 ms với dữ liệu mẫu & Lần chạy Postman ngày 21/10/2026: trung bình 14 ms, lớn nhất 108 ms\\
Băm mật khẩu, JWT, validate đầu vào & bcrypt trong \path|backend/src/auth/auth.service.ts|; \path|jwt-auth.guard.ts| chặn request thiếu token với mã 401; DTO dùng class-validator trả 400\\
Dùng được trên các trình duyệt web & \todo{trình duyệt và phiên bản đã thử}\\
\end{longtable}
\endgroup

\subsection{Kiến trúc và công nghệ}
\begingroup\small
\begin{longtable}{L{0.22\linewidth}L{0.70\linewidth}}
\toprule Thành phần & Công nghệ\\\midrule\endfirsthead
\toprule Thành phần & Công nghệ\\\midrule\endhead
\bottomrule\endfoot
Frontend & Next.js, TypeScript; giỏ hàng lưu bằng Zustand\\
Backend & NestJS, REST, class-validator, JWT\\
Cơ sở dữ liệu & PostgreSQL qua Prisma, đúng 5 thực thể Product, User, Order, OrderItem, Payment\\
Thanh toán & Cổng giả lập \path|backend/src/payment/mock-payment-gateway.service.ts|\\
Bàn giao & Docker Compose, README, collection Postman \path|docs/postman/BrewLite.postman_collection.json|\\
\end{longtable}
\endgroup
Backend tính lại giá từ bảng giá ở \path|backend/src/order/pricing.ts| và bỏ qua giá client gửi lên, nên khách không sửa được số tiền phải trả.

\subsection{Luồng đặt món và thanh toán}
\begin{enumerate}
\item Khách mở menu, chọn món, size và topping rồi thêm vào giỏ.
\item Khách đăng nhập nếu chưa có phiên, rồi bấm thanh toán.
\item Frontend gọi \texttt{POST} \path|/orders| để tạo đơn ở trạng thái Chờ thanh toán và nhận mã đơn.
\item Frontend gọi \texttt{POST} \path|/payments| với mã đơn, phương thức Ví hoặc Thẻ và khóa chống lặp \texttt{Idempotency-Key}, tức chuỗi ngẫu nhiên client gửi trong header để backend nhận ra request lặp lại.
\item Cổng chấp nhận thì đơn sang Đã thanh toán, web hiện biên lai. Cổng từ chối thì đơn sang Thanh toán lỗi, giỏ giữ nguyên và khách thử lại trên chính đơn đó.
\end{enumerate}

\section{Quy trình Agile Scrum}
\begingroup\small
\begin{longtable}{L{0.10\linewidth}L{0.42\linewidth}L{0.12\linewidth}L{0.22\linewidth}}
\toprule Sprint & Mục tiêu (mục 2.3 của đề) & Task & Thời gian\\\midrule\endfirsthead
\toprule Sprint & Mục tiêu (mục 2.3 của đề) & Task & Thời gian\\\midrule\endhead
\bottomrule\endfoot
1 & Nền tảng và hiển thị menu & 1, 2, 3, 4 & \todo{ngày}\\
2 & Giỏ hàng, đặt đơn, đăng nhập & 5, 6, 7 & \todo{ngày}\\
3 & Thanh toán, bàn giao, nghiệp vụ nâng cao & 8, 9, 10 & \todo{ngày}\\
\end{longtable}
\endgroup
Product Backlog gồm 10 user story ở \path|docs/product-backlog.md|, mỗi story ứng với một task của đề. Một task chỉ tính xong khi đạt Definition of Done ở mục 10 của đề.

Product Owner: \todo{tên}. Scrum Master: \todo{tên}. Lịch Sprint Planning, Daily Scrum và Sprint Review: \todo{ngày và người dự}.

Burndown dưới đây tính cho cả 10 task. Một task tính là xong vào ngày tác giả tạo commit của thành viên phụ trách, theo bảng ở mục 4. Đường lý tưởng giảm đều từ 10 task ngày 30/09/2026, ngày có commit đầu tiên của repo, về 0 ngày 23/10/2026, ngày Review 1.

\begin{figure}[H]\centering
\khunganh{0.9\linewidth}{burndown.png}
\caption{Burndown 10 task dựng từ ngày commit, sinh bằng \texttt{burndown.py}}\label{fig:burndown}
\end{figure}

\begingroup\small
\begin{longtable}{L{0.12\linewidth}L{0.30\linewidth}L{0.20\linewidth}L{0.12\linewidth}L{0.12\linewidth}}
\toprule Ngày & Task có commit & Commit & Thực tế & Lý tưởng\\\midrule\endfirsthead
\toprule Ngày & Task có commit & Commit & Thực tế & Lý tưởng\\\midrule\endhead
\bottomrule\endfoot
30/09 & Không có, repo được khởi tạo & \texttt{0d022c0} & 10 & 10\\
07/10 & Task 3, 5, 8, 10 & \texttt{9ea0933}, \texttt{bd1adc3} & 6 & 7,0\\
13/10 & Task 1, 4, 6, 7 & \texttt{44e3784}, \texttt{32831dc} & 2 & 4,3\\
21/10 & Task 2, 9 & \texttt{a358930}, \texttt{ed2e4ee} & 0 & 0,9\\
23/10 & Không có, ngày Review 1 & Không có & 0 & 0\\
\end{longtable}
\endgroup
Những ngày không có trong bảng giữ nguyên số task của ngày liền trước. Theo cách đọc ở mục 2.4 của đề, đường thực tế nằm trên đường lý tưởng là chậm tiến độ. Nhóm chậm trong ba khoảng: 01/10 đến 06/10 khi chưa task nào có commit, 10/10 đến 12/10 trước đợt commit ngày 13/10, và 19/10 đến 20/10 khi Task 2 và Task 9 chưa có commit. Các ngày còn lại nhóm đi trước đường lý tưởng. Không có ngày bắt đầu và kết thúc từng Sprint nên chưa tách được burndown riêng cho mỗi Sprint.

\section{Phân công và commit}
Thuận commit dưới bốn tên Git: Thuan-Test, Trần Vĩnh Thuận, Thuan-Lap và Thuan-AFK666. Toàn bộ dữ liệu Git nằm ở \path|git-evidence.txt|.

\begingroup\small
\begin{longtable}{L{0.16\linewidth}L{0.18\linewidth}L{0.24\linewidth}L{0.30\linewidth}}
\toprule Thành viên & Commit, tác giả & Nhánh & Nội dung\\\midrule\endfirsthead
\toprule Thành viên & Commit, tác giả & Nhánh & Nội dung\\\midrule\endhead
\bottomrule\endfoot
Bảo Tín, Task 1 & \texttt{44e3784}, Saturlen & \texttt{main} & README, Docker Compose, Dockerfile, mã backend\\\addlinespace
Phú, Task 2 & \texttt{a358930}, \texttt{5d92e1f}, phu0111 & \texttt{a1task2-9} & Đặc tả API danh sách sản phẩm\\\addlinespace
Bảo Thiện, Task 3 & \texttt{9ea0933}, baothien-s & \texttt{thien35}, merge \texttt{862bf65} & Trang Menu, gọi API, trạng thái tải và rỗng\\\addlinespace
Lợi, Task 4 & \texttt{32831dc}, DuVinhLoi & \texttt{feat/task4-task7}, merge \texttt{eaf33f8} & Chi tiết món, trạng thái tải và lỗi, test\\\addlinespace
Bảo Thiện, Task 5 & \texttt{9ea0933}, baothien-s & \texttt{thien35}, merge \texttt{862bf65} & Giỏ hàng, tổng tiền, badge\\\addlinespace
Bảo Tín, Task 6 & \texttt{44e3784}, Saturlen & \texttt{main} & Controller, DTO và bảng giá của API tạo đơn\\\addlinespace
Lợi, Task 7 & \texttt{32831dc}, DuVinhLoi & \texttt{feat/task4-task7}, merge \texttt{eaf33f8} & JWT, chuẩn hóa email, test xác thực\\\addlinespace
Thuận, Task 8 & \texttt{bd1adc3}, Thuan-Test & \path|feature/task-8-10-payment-order|, merge \texttt{03cb553} & Thanh toán qua cổng giả lập, chuyển trạng thái sau thanh toán\\\addlinespace
Phú, Task 9 & \texttt{ed2e4ee}, phu0111 & \texttt{a1task2-9} & Đặc tả API xác nhận đơn và lịch sử đơn\\\addlinespace
Thuận, Task 10 & \texttt{bd1adc3}, Thuan-Test & \path|feature/task-8-10-payment-order|, merge \texttt{03cb553} & Máy trạng thái, idempotency, tồn kho, điểm thưởng\\
\end{longtable}
\endgroup

\section{Minh chứng cho 10 task}
Mỗi task gồm tiêu chí của đề, các ca kiểm thử trong collection Postman, file mã nguồn, test tự động và ảnh chụp. Request trong collection đặt tên theo mã ca, ví dụ \texttt{T06.5} là ca thứ 5 của Task 6. Chi tiết từng ca nằm ở \path|docs/minh-chung-10-task.md|.

Ảnh chụp đặt vào thư mục \path|figures/| cạnh \path|main.tex|, đúng tên file ghi trong khung. Khung ghi ``Chưa có ảnh'' là ảnh nhóm chưa chụp.

\subsection{Task 1 — Khởi tạo dự án và cấu trúc}
\textbf{Tiêu chí của đề.} Repo có thư mục \path|frontend| chạy Next.js và \path|backend| chạy NestJS, chạy được trang chào, có README, Git và \path|.env.example|.

\begingroup\small
\begin{bangca}
T01.1 & \texttt{GET} \path|{{webUrl}}/| & 200, HTML có chữ ``BrewLite''\\
T01.2 & \texttt{GET} \path|{{baseUrl}}/products| & 200, JSON\\
\end{bangca}
\endgroup

\textbf{Mã nguồn.} \path|README.md|, \path|.env.example|, \path|frontend/package.json|, \path|backend/package.json|, \path|docker-compose.yml|.

\textbf{Test.} \path|backend/test/bootstrap-env.spec.ts| kiểm backend dừng khi thiếu \path|DATABASE_URL|; \path|backend/test/task47-auth.spec.ts| kiểm backend dừng khi thiếu \path|JWT_SECRET|.

\begin{figure}[H]\centering
\begin{minipage}[t]{0.48\linewidth}\centering\khunganh{\linewidth}{task01-tree.png}\par\footnotesize (a) Cây thư mục repo\end{minipage}\hfill
\begin{minipage}[t]{0.48\linewidth}\centering\khunganh{\linewidth}{task01-gitlog.png}\par\footnotesize (b) Kết quả \texttt{git log -{}-oneline}\end{minipage}\par\medskip
\begin{minipage}[t]{0.48\linewidth}\centering\khunganh{\linewidth}{task01-runner.png}\par\footnotesize (c) Postman Runner thư mục Task 01\end{minipage}
\caption{Minh chứng Task 1}\label{fig:task01}
\end{figure}

\subsection{Task 2 — API danh sách sản phẩm}
\textbf{Tiêu chí của đề.} \texttt{GET} \path|/products| trả JSON menu gồm id, tên, giá, ảnh.

\begingroup\small
\begin{bangca}
T02.1 & \texttt{GET} \path|/products| & Mảng món, mỗi món có \texttt{id}, \texttt{name}, \texttt{price}, \texttt{imageUrl}; không có món \texttt{stock = 0}\\
T02.2 & \texttt{GET} \path|/products/:id| với id có thật & Đúng món đó\\
T02.3 & \texttt{GET} \path|/products/:id| với id không có & 404\\
\end{bangca}
\endgroup

\textbf{Mã nguồn.} \path|backend/src/product/product.controller.ts|, dữ liệu mẫu \path|backend/prisma/seed.ts|.

\textbf{Test.} \path|backend/test/order-product.spec.ts|, ca ``product list excludes stock-0'' và ``product detail unknown → NotFound''.

\begin{figure}[H]\centering
\khunganh{0.8\linewidth}{task02-T02.1.png}
\caption{Postman ca T02.1: tab Body và Test Results}\label{fig:task02}
\end{figure}

\subsection{Task 3 — Trang Menu}
\textbf{Tiêu chí của đề.} Next.js gọi \path|/products|, hiện lưới sản phẩm, có trạng thái đang tải và danh sách rỗng.

\begingroup\small
\begin{bangca}
T03.1 & \texttt{GET} \path|{{webUrl}}/| & 200\\
T03.2 & Tải ảnh theo \texttt{imageUrl} của một món & 200, kiểu \texttt{image/svg+xml}\\
\end{bangca}
\endgroup

Postman chỉ kiểm được phía server. Phần hiển thị dựa vào test vitest và ảnh chụp trình duyệt.

\textbf{Mã nguồn.} \path|frontend/src/components/ProductGrid.tsx|, hàm \texttt{getProducts} trong \path|frontend/src/lib/api.ts|.

\textbf{Test.} \path|frontend/src/components/ProductGrid.test.ts| kiểm ba trạng thái đang tải, rỗng và lỗi kèm nút ``Thử lại''; \path|frontend/src/lib/api.test.ts|.

\begin{figure}[H]\centering
\begin{minipage}[t]{0.48\linewidth}\centering\khunganh{\linewidth}{task03-menu.png}\par\footnotesize (a) Lưới món tại \path|http://localhost:3001|\end{minipage}\hfill
\begin{minipage}[t]{0.48\linewidth}\centering\khunganh{\linewidth}{task03-loading.png}\par\footnotesize (b) Đang tải, mạng giả lập Slow 3G trong DevTools\end{minipage}\par\medskip
\begin{minipage}[t]{0.48\linewidth}\centering\khunganh{\linewidth}{task03-error.png}\par\footnotesize (c) Lỗi khi tắt backend bằng \texttt{docker compose stop backend}\end{minipage}
\caption{Minh chứng Task 3}\label{fig:task03}
\end{figure}

\subsection{Task 4 — Chi tiết và tùy chọn sản phẩm}
\textbf{Tiêu chí của đề.} Chọn size S, M, L và topping; giá đổi theo tùy chọn.

Size M cộng 5.000 đồng, size L cộng 10.000 đồng. Topping Trân châu giá 5.000 đồng, Kem cheese giá 8.000 đồng. Frontend và backend dùng cùng bảng giá này.

\begingroup\small
\begin{bangca}
T04.1 & Đặt size L, topping Trân châu và Kem cheese & \texttt{unitPrice} = giá gốc + 10.000 + 13.000\\
T04.2 & Đặt size S, không topping & \texttt{unitPrice} = giá gốc\\
T04.3 & Đặt với topping không có trong bảng giá & 400 \texttt{Unknown topping}\\
\end{bangca}
\endgroup

\textbf{Mã nguồn.} \path|frontend/src/components/ProductDetailModal.tsx|, \path|frontend/src/lib/pricing.ts|, \path|backend/src/order/pricing.ts|.

\textbf{Test.} \path|frontend/src/lib/pricing.test.ts|, \path|frontend/src/components/ProductDetailModal.test.ts|, \path|frontend/src/lib/product-detail.test.ts|; \path|backend/test/pricing-table.spec.ts| kiểm bảng giá hai phía trùng nhau.

\begin{figure}[H]\centering
\begin{minipage}[t]{0.48\linewidth}\centering\khunganh{\linewidth}{task04-before.png}\par\footnotesize (a) Trước khi đổi size và topping\end{minipage}\hfill
\begin{minipage}[t]{0.48\linewidth}\centering\khunganh{\linewidth}{task04-after.png}\par\footnotesize (b) Sau khi đổi, giá tăng tương ứng\end{minipage}
\caption{Modal chi tiết món}\label{fig:task04}
\end{figure}

\subsection{Task 5 — Giỏ hàng}
\textbf{Tiêu chí của đề.} Thêm, sửa số lượng, xóa; lưu state; tự tính tổng; badge số lượng.

Giỏ giới hạn 20 dòng và 1 đến 99 ly mỗi dòng; backend kiểm lại cùng giới hạn khi tạo đơn.

\begingroup\small
\begin{bangca}
T05.1 & Tạo đơn từ giỏ 2 dòng & \texttt{total} = tổng \texttt{quantity × unitPrice}\\
T05.2 & Tạo đơn từ giỏ 21 dòng & 400\\
T05.3 & Tạo đơn có một dòng 100 ly & 400\\
\end{bangca}
\endgroup

\textbf{Mã nguồn.} Store Zustand \path|frontend/src/store/cart.ts|, \path|frontend/src/components/cart/CartView.tsx|, badge trong \path|frontend/src/components/HeaderView.tsx|.

\textbf{Test.} \path|frontend/src/store/cart.test.ts|, \path|frontend/src/lib/cart-limits.test.ts|, \path|frontend/src/components/cart/CartView.test.ts|, \path|frontend/src/components/QtyStepper.test.ts|, \path|frontend/src/components/HeaderView.test.ts|.

\begin{figure}[H]\centering
\begin{minipage}[t]{0.48\linewidth}\centering\khunganh{\linewidth}{task05-cart.png}\par\footnotesize (a) Giỏ có 2 món và tổng tiền\end{minipage}\hfill
\begin{minipage}[t]{0.48\linewidth}\centering\khunganh{\linewidth}{task05-badge.png}\par\footnotesize (b) Badge số lượng trên header\end{minipage}\par\medskip
\begin{minipage}[t]{0.48\linewidth}\centering\khunganh{\linewidth}{task05-reload.png}\par\footnotesize (c) Giỏ còn nguyên sau khi tải lại trang\end{minipage}
\caption{Minh chứng Task 5}\label{fig:task05}
\end{figure}

\subsection{Task 6 — API tạo đơn hàng}
\textbf{Tiêu chí của đề.} \texttt{POST} \path|/orders| nhận giỏ hàng, validate bằng class-validator, lưu đơn ở trạng thái Chờ thanh toán và trả mã đơn.

\begingroup\small
\begin{bangca}
T06.1 & Gửi không kèm token & 401\\
T06.2 & Giỏ rỗng & 400\\
T06.3 & Size \texttt{XL} & 400\\
T06.4 & Món không có trong menu & 404\\
T06.5 & Đơn hợp lệ, cố tình gửi kèm \texttt{unitPrice: 1} & 201, \texttt{status} = \texttt{PENDING}, có \texttt{id}; giá do server tính, bỏ qua giá 1 đồng\\
T06.6 & Gửi bằng token nhân viên quầy & 403\\
\end{bangca}
\endgroup

\textbf{Mã nguồn.} \path|backend/src/order/order.controller.ts|, \path|backend/src/order/dto/create-order.dto.ts|, \path|backend/src/order/pricing.ts|.

\textbf{Test.} \path|backend/test/order-product.spec.ts|, \path|backend/test/pricing-table.spec.ts|, \path|backend/test/customer-only.spec.ts|.

\begin{figure}[H]\centering
\begin{minipage}[t]{0.48\linewidth}\centering\khunganh{\linewidth}{task06-T06.5.png}\par\footnotesize (a) Ca T06.5: Body có \texttt{id} và \texttt{status}\end{minipage}\hfill
\begin{minipage}[t]{0.48\linewidth}\centering\khunganh{\linewidth}{task06-T06.2.png}\par\footnotesize (b) Ca T06.2: 400 kèm thông báo validate\end{minipage}
\caption{Minh chứng Task 6}\label{fig:task06}
\end{figure}

\subsection{Task 7 — Đăng ký và đăng nhập bằng JWT}
\textbf{Tiêu chí của đề.} \texttt{POST} \path|/auth/register| và \texttt{POST} \path|/auth/login|; băm mật khẩu bằng bcrypt; phát JWT; guard bảo vệ route đặt đơn.

\begingroup\small
\begin{bangca}
T07.1 & Đăng ký email mới & 201, body không có \texttt{passwordHash}\\
T07.2 & Đăng ký lại cùng email viết hoa & 409\\
T07.3 & Email sai định dạng, mật khẩu quá ngắn & 400\\
T07.4 & Đăng nhập đúng & JWT đủ 3 phần, \texttt{role} = \texttt{CUSTOMER}\\
T07.5 & Đăng nhập sai mật khẩu & 401 \texttt{Invalid credentials}\\
T07.6 & Đăng nhập email không tồn tại & 401 \texttt{Invalid credentials}\\
T07.7 & Gọi API bằng token giả & 401\\
\end{bangca}
\endgroup

Ca sai mật khẩu và ca email không tồn tại trả cùng một thông báo, nên người ngoài không dò được email nào đã đăng ký. Guard của route đặt đơn được kiểm ở ca T06.1.

\textbf{Mã nguồn.} \path|backend/src/auth/auth.controller.ts|, \path|backend/src/auth/auth.service.ts|, \path|backend/src/auth/jwt-auth.guard.ts|, \path|frontend/src/app/login/page.tsx|.

\textbf{Test.} \path|backend/test/auth.spec.ts|, \path|backend/test/task47-auth.spec.ts|, \path|backend/test/roles.spec.ts|, \path|frontend/src/lib/login.test.ts|, \path|frontend/src/lib/session.test.ts|, \path|frontend/src/components/LoginView.test.ts|.

\begin{figure}[H]\centering
\begin{minipage}[t]{0.48\linewidth}\centering\khunganh{\linewidth}{task07-T07.4.png}\par\footnotesize (a) Ca T07.4: có \texttt{accessToken}\end{minipage}\hfill
\begin{minipage}[t]{0.48\linewidth}\centering\khunganh{\linewidth}{task07-T07.5.png}\par\footnotesize (b) Ca T07.5: 401\end{minipage}\par\medskip
\begin{minipage}[t]{0.8\linewidth}\centering\khunganh{\linewidth}{task07-db-hash.png}\par\footnotesize (c) Bảng \texttt{User}: \texttt{passwordHash} bắt đầu bằng \texttt{\$2a\$}, tức chuỗi băm bcrypt\end{minipage}
\caption{Minh chứng Task 7}\label{fig:task07}
\end{figure}

Lệnh lấy ảnh (c):
\begin{verbatim}
docker compose exec db psql -U brewlite \
  -c 'select email, left("passwordHash", 7) from "User" limit 3;'
\end{verbatim}

\subsection{Task 8 — Thanh toán không tiền mặt}
\textbf{Tiêu chí của đề.} Chọn Ví hoặc Thẻ, gọi \texttt{POST} \path|/payments|. Thanh toán thành công thì đơn sang Đã thanh toán; lỗi thì đơn sang Thanh toán lỗi và giỏ giữ nguyên.

Để tái hiện luồng thanh toán lỗi, nhóm dựng stack với biến môi trường \path|MOCK_GATEWAY_DECLINE_METHOD=CARD|. Biến này buộc cổng giả lập từ chối mọi giao dịch bằng Thẻ, còn Ví vẫn thành công.

\begingroup\small
\begin{bangca}
T08.1 & Trả bằng Ví & 201, \texttt{PAID}, có \texttt{paymentId}\\
T08.2 & Đọc lại đơn vừa trả & \texttt{status} = \texttt{PAID}\\
T08.3--T08.5 & Trả bằng Thẻ, cổng giả lập từ chối & 400 \texttt{Payment failed}; đơn sang \texttt{PAYMENT\_FAILED}\\
T08.6 & Trả lại bằng Ví trên đơn vừa lỗi & 201, \texttt{PAID}, cùng mã đơn cũ\\
T08.7 & Thiếu header \texttt{Idempotency-Key} & 400\\
T08.8 & Phương thức \texttt{CASH} & 400\\
\end{bangca}
\endgroup

\textbf{Mã nguồn.} \path|backend/src/payment/payment.controller.ts|, \path|backend/src/payment/payment.service.ts|, \path|backend/src/payment/mock-payment-gateway.service.ts|, \path|frontend/src/components/PaymentForm.tsx|.

\textbf{Test.} \path|backend/test/payment-failure.spec.ts|, \path|backend/test/order-payment.integration.spec.ts|, \path|frontend/src/components/PaymentForm.test.ts|, \path|frontend/src/lib/payment/api-gateway.test.ts|, \path|frontend/src/lib/payment/checkout-session.test.ts|.

\begin{figure}[H]\centering
\begin{minipage}[t]{0.48\linewidth}\centering\khunganh{\linewidth}{task08-T08.4.png}\par\footnotesize (a) Ca T08.4: 400 \texttt{Payment failed}\end{minipage}\hfill
\begin{minipage}[t]{0.48\linewidth}\centering\khunganh{\linewidth}{task08-T08.5.png}\par\footnotesize (b) Ca T08.5: đơn ở \texttt{PAYMENT\_FAILED}\end{minipage}\par\medskip
\begin{minipage}[t]{0.48\linewidth}\centering\khunganh{\linewidth}{task08-web-error.png}\par\footnotesize (c) Màn thanh toán báo lỗi khi chọn Thẻ, giỏ còn nguyên\end{minipage}
\caption{Minh chứng Task 8}\label{fig:task08}
\end{figure}

\subsection{Task 9 — Xác nhận, lịch sử đơn và bàn giao}
\textbf{Tiêu chí của đề.} Màn xác nhận có mã đơn và trạng thái; \texttt{GET} \path|/orders/me| trả lịch sử đơn; Docker Compose chạy đủ frontend, backend và cơ sở dữ liệu; có README; demo được trọn luồng.

\begingroup\small
\begin{bangca}
T09.1 & \texttt{GET} \path|/orders/:id| đơn của mình & Có \texttt{id}, \texttt{status}, \texttt{total}, \texttt{items}\\
T09.2 & \texttt{GET} \path|/orders/me| & Có các đơn vừa tạo, mới nhất trước, mỗi dòng có \texttt{productName}\\
T09.3 & Khách B gọi \texttt{GET} \path|/orders/me| & Không thấy đơn của khách A\\
T09.4 & Khách B đọc đơn của khách A & 403\\
T09.5 & Mã đơn không tồn tại & 404\\
T09.6 & \texttt{GET} \path|{{webUrl}}/history| & 200\\
\end{bangca}
\endgroup

\textbf{Mã nguồn.} \path|backend/src/order/order.controller.ts|, \path|frontend/src/app/history/page.tsx|, \path|frontend/src/components/PaymentSuccess.tsx|, \path|docker-compose.yml|, \path|README.md|.

\textbf{Test.} \path|backend/test/order-product.spec.ts|, ca ``listMine returns newest-first''; \path|backend/test/compose-env.spec.ts|; \path|frontend/src/components/HistoryView.test.ts|, \path|frontend/src/components/PaymentSuccess.test.ts|, \path|frontend/src/lib/history.test.ts|, \path|frontend/src/lib/use-order-poll.test.ts|.

\begin{figure}[H]\centering
\begin{minipage}[t]{0.48\linewidth}\centering\khunganh{\linewidth}{task09-compose-ps.png}\par\footnotesize (a) \texttt{docker compose ps} có \texttt{db}, \texttt{backend}, \texttt{frontend}\end{minipage}\hfill
\begin{minipage}[t]{0.48\linewidth}\centering\khunganh{\linewidth}{task09-receipt.png}\par\footnotesize (b) Biên lai có mã đơn, trạng thái Đã thanh toán\end{minipage}\par\medskip
\begin{minipage}[t]{0.48\linewidth}\centering\khunganh{\linewidth}{task09-history.png}\par\footnotesize (c) Trang \path|/history|\end{minipage}
\caption{Minh chứng Task 9}\label{fig:task09}
\end{figure}

\subsection{Task 10 — Nghiệp vụ backend}
\textbf{Tiêu chí của đề.} Mục 9 của đề yêu cầu đủ bốn phần: máy trạng thái đơn hàng, thanh toán idempotent, kiểm soát tồn kho khi đặt đồng thời, khuyến mãi và điểm thưởng. Test phải chứng minh ba điều: chặn được chuyển trạng thái sai; gửi hai lần cùng \texttt{Idempotency-Key} chỉ tạo một đơn; đặt đồng thời không vượt tồn kho.

\textbf{(a) Máy trạng thái đơn hàng.} Hàm \texttt{assertTransition} trong \path|backend/src/state-machine/order-state.machine.ts| chặn mọi chuyển tiếp không có trong bảng dưới với mã 400. Hoàn tất và Đã hủy là trạng thái cuối. Quyền được kiểm riêng trước khi kiểm chuyển tiếp: khách gọi API tiến đơn bị chặn với mã 403 dù chuyển tiếp đó hợp lệ.

\begingroup\small
\begin{longtable}{L{0.22\linewidth}L{0.22\linewidth}L{0.48\linewidth}}
\toprule Từ & Sang & Bên thực hiện\\\midrule\endfirsthead
\toprule Từ & Sang & Bên thực hiện\\\midrule\endhead
\bottomrule\endfoot
\texttt{PENDING} & \texttt{PAID} & Backend, khi cổng giả lập chấp nhận giao dịch\\
\texttt{PENDING} & \texttt{PAYMENT\_FAILED} & Backend, khi cổng giả lập từ chối\\
\texttt{PAYMENT\_FAILED} & \texttt{PENDING} & Backend, khi khách thanh toán lại đơn đó\\
\texttt{PAID} & \texttt{PREPARING} & Barista hoặc Admin\\
\texttt{PREPARING} & \texttt{READY} & Barista hoặc Admin\\
\texttt{READY} & \texttt{COMPLETED} & Barista hoặc Admin\\
\texttt{PENDING} & \texttt{CANCELLED} & Khách chủ đơn\\
\texttt{PAYMENT\_FAILED} & \texttt{CANCELLED} & Khách chủ đơn\\
\texttt{PAID} & \texttt{CANCELLED} & Khách chủ đơn, hoặc Barista và Admin khi hết nguyên liệu\\
\end{longtable}
\endgroup
Giao diện hiển thị các trạng thái theo nhãn trong \path|frontend/src/lib/order-status.ts|: \texttt{PENDING} là Chờ thanh toán, \texttt{PAID} là Đã thanh toán, \texttt{PREPARING} là Đang pha, \texttt{READY} là Sẵn sàng, \texttt{COMPLETED} là Hoàn tất, \texttt{PAYMENT\_FAILED} là Thanh toán lỗi, \texttt{CANCELLED} là Đã hủy.

\begingroup\small
\begin{bangca}
T10a.1 & Khách tự tiến đơn của mình & 403 \texttt{Only staff can advance orders}\\
T10a.2 & Quầy chuyển thẳng \texttt{PAID} sang \texttt{READY} & 400 \texttt{Invalid transition PAID -> READY}\\
T10a.3 & Quầy tiến đơn sang \texttt{PREPARING} & 200\\
T10a.4 & Khách hủy đơn đang \texttt{PREPARING} & 400\\
T10a.5, T10a.6 & Quầy tiến đơn sang \texttt{READY}, rồi \texttt{COMPLETED} & 200 mỗi bước\\
T10a.7 & Hủy đơn đã \texttt{COMPLETED} & 400 \texttt{Invalid transition COMPLETED -> CANCELLED}\\
T10a.8 & Đọc lại đơn & Vẫn \texttt{COMPLETED}\\
T10a.9 & Quầy gọi \texttt{GET} \path|/orders?status=PAID| & Có đơn vừa thanh toán\\
T10a.10 & Quầy hủy đơn \texttt{PAID} & 200, \texttt{CANCELLED}\\
\end{bangca}
\endgroup

\textbf{Mã nguồn.} \path|backend/src/state-machine/order-state.machine.ts|; hàm \texttt{advanceOrder} và \texttt{cancelOrderAs} trong \path|backend/src/payment/payment.service.ts|.

\textbf{Test.} \path|backend/src/state-machine/order-state.machine.spec.ts|, \path|backend/test/order-advance.spec.ts|, \path|backend/test/barista-queue.spec.ts|, \path|frontend/src/lib/order-status.test.ts|.

\textbf{(b) Thanh toán idempotent.} Backend lưu \texttt{Idempotency-Key} với ràng buộc duy nhất ở bảng Order và bảng Payment. Request lặp lại cùng khóa nhận lại kết quả cũ với mã 200, không tạo đơn mới, không trừ tiền, không trừ kho và không cộng điểm thêm lần nữa.

\begingroup\small
\begin{bangca}
T10b.2, T10b.3 & \texttt{POST} \path|/orders| hai lần cùng khóa & Lần 1 trả 201; lần 2 trả 200, cùng \texttt{id}, \texttt{isCached: true}\\
T10b.5, T10b.6 & \texttt{POST} \path|/payments| hai lần cùng khóa & Lần 2 trả 200, cùng \texttt{paymentId}\\
T10b.7 & Đọc điểm thưởng & Chỉ cộng một lần\\
T10b.8 & Đọc tồn kho & Chỉ trừ một lần, đúng 2 ly\\
T10b.9 & Khách khác dùng lại khóa đó & 422\\
\end{bangca}
\endgroup

\textbf{Test.} \path|backend/test/order-payment.integration.spec.ts|, ca ``(b) same Idempotency-Key charges once''.

\textbf{(c) Tồn kho khi đặt đồng thời.} Kho bị trừ lúc thanh toán, trong cùng transaction với việc chuyển đơn sang Đã thanh toán, ghi điểm thưởng và lưu biên lai. Mỗi lần trừ kho so cột \texttt{Product.version}: bản ghi đã bị request khác sửa thì lần ghi thất bại và backend đọc lại, thử tối đa 3 lần. Cách khóa này gọi là optimistic locking. Hết 3 lần vẫn không trừ được thì transaction rollback, backend hoàn tiền và đơn sang Thanh toán lỗi.

\begingroup\small
\begin{bangca}
T10c.1 & Admin ghi kho với \texttt{version} cũ & 409 \texttt{Stock changed, reload}\\
T10c.2 & Admin trừ quá số còn lại & 400 \texttt{Stock cannot be negative}\\
T10c.3 & Khách gọi API kho & 403\\
\end{bangca}
\endgroup

Postman Runner gửi request lần lượt nên không tạo được tải đồng thời. Trường hợp nhiều khách đặt cùng lúc được chứng minh bằng test ``(c) concurrent stock never negative'' trong \path|backend/test/order-payment.integration.spec.ts| và \path|backend/test/pg-payment.integration.spec.ts|. Test thứ hai chạy trên PostgreSQL thật và chỉ chạy khi có biến \path|DATABASE_URL|; lần chạy ngày 21/10/2026 không đặt biến này nên 2 ca của nó bị bỏ qua.

\textbf{(d) Khuyến mãi và điểm thưởng.} Đề không quy định mã và tỉ lệ, nhóm chọn mã \texttt{BREWLITE10} giảm 10\%:
\[
\text{giảm giá}=\lfloor \text{tổng}\times0{,}1\rfloor,\qquad
\text{số tiền trả}=\text{tổng}-\text{giảm giá},\qquad
\text{điểm}=\left\lfloor \frac{\text{số tiền trả}\times0{,}1}{1000}\right\rfloor.
\]
Điểm chỉ cộng sau khi đơn sang Đã thanh toán. Lần chạy ngày 21/10/2026: đơn 58.000 đồng được giảm 5.800 đồng, khách trả 52.200 đồng và được cộng 5 điểm.

\begingroup\small
\begin{bangca}
T10d.3 & Trả kèm mã \texttt{BREWLITE10} & \texttt{discount}, \texttt{amount}, \texttt{loyaltyPointsEarned} đúng công thức\\
T10d.4 & \texttt{GET} \path|/users/me| & Thấy điểm vừa cộng\\
T10d.6 & Trả kèm mã không hợp lệ \texttt{SAIMA} & \texttt{discount} = 0\\
\end{bangca}
\endgroup

\textbf{Test.} \path|backend/test/order-payment.integration.spec.ts|, các ca của mã \texttt{BREWLITE10}.

\begin{figure}[H]\centering
\begin{minipage}[t]{0.48\linewidth}\centering\khunganh{\linewidth}{task10-runner.png}\par\footnotesize (a) Postman Runner thư mục Task 10, đủ 4 thư mục con\end{minipage}\hfill
\begin{minipage}[t]{0.48\linewidth}\centering\khunganh{\linewidth}{task10-jest.png}\par\footnotesize (b) \texttt{npx jest order-state order-advance order-payment} chạy xanh\end{minipage}\par\medskip
\begin{minipage}[t]{0.48\linewidth}\centering\khunganh{\linewidth}{task10-history.png}\par\footnotesize (c) Trang \path|/history| tự đổi trạng thái khi quầy tiến đơn\end{minipage}
\caption{Minh chứng Task 10}\label{fig:task10}
\end{figure}

\section{Kiểm thử, Definition of Done và bàn giao}
\subsection{Kết quả kiểm thử ngày 21/10/2026}
\begingroup\small
\begin{longtable}{L{0.30\linewidth}L{0.62\linewidth}}
\toprule Bộ kiểm thử & Kết quả\\\midrule\endfirsthead
\toprule Bộ kiểm thử & Kết quả\\\midrule\endhead
\bottomrule\endfoot
Postman, 74 request & 137 assertion đạt, 0 lỗi; phản hồi trung bình 14 ms, lớn nhất 108 ms\\
Backend, \texttt{npx jest} & 170 test đạt, 3 test bỏ qua, 173 tổng\\
Frontend, \texttt{npx vitest run} & 733 test đạt trong 46 file\\
\end{longtable}
\endgroup
Ba test backend bị bỏ qua đều cần PostgreSQL thật qua biến \path|DATABASE_URL|: 2 ca idempotency và tồn kho đồng thời trong \path|pg-payment.integration.spec.ts|, 1 ca dọn dữ liệu thử trong \path|pg-smoke-clean.integration.spec.ts|. Log Postman lưu ở \path|docs/postman/results/2026-10-21-run.txt|.

\subsection{Definition of Done}
\begingroup\small
\begin{longtable}{L{0.28\linewidth}L{0.36\linewidth}L{0.28\linewidth}}
\toprule Điều kiện (mục 10 của đề) & Cách xác nhận & Trạng thái\\\midrule\endfirsthead
\toprule Điều kiện (mục 10 của đề) & Cách xác nhận & Trạng thái\\\midrule\endhead
\bottomrule\endfoot
Code chạy, không lỗi build & Dựng stack bằng \texttt{docker compose up -d -{}-build} & Dựng được ngày 21/10/2026 để chạy Postman\\
Đúng tiêu chí chấp nhận của task & Ca Postman và test ở mục 5 & 10 task có ca kiểm thử; còn thiếu ảnh\\
Có commit rõ ràng, được review trong nhóm & Commit ở mục 4 & Có commit; \todo{người review từng task}\\
API validate đầu vào, UI xử lý đang tải và lỗi & Ca 400 của Task 6, 7, 8; test \path|ProductGrid.test.ts| & Đạt theo lần chạy 21/10/2026\\
README có hướng dẫn chạy & Người ngoài nhóm làm theo README trên máy mới & \todo{kết quả thử trên máy mới}\\
\end{longtable}
\endgroup

\subsection{Kịch bản demo bàn giao}
\begin{enumerate}
\item Dựng ba dịch vụ theo README, chụp \texttt{docker compose ps}.
\item Đăng ký, đăng nhập, xem menu, đổi size và topping, xem giá đổi.
\item Thêm hai món, sửa số lượng, xóa một món, xem badge và tổng tiền.
\item Thanh toán bằng Ví, xem biên lai có mã đơn và trang lịch sử.
\item Thanh toán bằng Thẻ khi cổng giả lập từ chối, xem giỏ còn nguyên, rồi thanh toán lại bằng Ví trên cùng đơn.
\item Chạy thư mục Task 10 trong Postman: tiến đơn bằng tài khoản quầy, gửi lặp cùng khóa, ghi kho với \texttt{version} cũ.
\item Chạy test tồn kho đồng thời trên PostgreSQL thật.
\item Bàn giao mã nguồn, \path|.env.example|, README, collection Postman và log kiểm thử; ghi phản hồi của khách hàng vào Product Backlog.
\end{enumerate}

\subsection{Sprint Retrospective}
Điều làm tốt: \todo{ý kiến của nhóm}. Điều cần cải thiện: \todo{ý kiến của nhóm}. Hành động cho Sprint sau: \todo{ý kiến của nhóm}.

\section{Nguồn và cách chạy lại}
\begin{itemize}
\item \path|docs/BrewLite-CongNghePhanMem-fixed.md|: đề, tiêu chí 10 task và Definition of Done.
\item \path|docs/minh-chung-10-task.md|: chi tiết từng ca Postman, test và ảnh cần chụp.
\item \path|docs/product-backlog.md|: 10 user story và tiêu chí chấp nhận.
\item \path|docs/postman/BrewLite.postman_collection.json| và \path|docs/postman/results/2026-10-21-run.txt|: collection và log Postman.
\item \path|git-evidence.txt|: commit của từng thành viên.
\end{itemize}
Chạy lại kiểm thử từ thư mục gốc của repo:
\begin{verbatim}
MOCK_GATEWAY_DECLINE_METHOD=CARD docker compose up -d --build
postman collection run docs/postman/BrewLite.postman_collection.json
(cd backend && npx jest)
(cd frontend && npx vitest run)
\end{verbatim}
\end{document}
